\section{General Framework}
\label{sec:general-framework}

Our general framework includes: K players, each parametrized by their own $\theta_k$, a collection of tasks/games that the players have to perform, a communication protocol $V$ that enables the players to communicate with each other in order to solve the task and payoffs assigned to the players as a deterministic function of a well-defined goal. We focus for now on a specific problem: \textit{referential games} with the following structure:

\begin{enumerate}
\item There is a set of images represented by vectors $\lbrace i_1, \dots, i_N \rbrace$, two images are drawn at random from this set, call them $(i_L, i_R)$, one of them is chosen to be the ``target'' $t \in \lbrace L, R \rbrace$
\item There are two players, a sender and a receiver each seeing the images - the sender receives input $\theta_S (i_L, i_R, t)$
\item There is a \textit{vocabulary} $V$ of size $K$ and the sender chooses one symbol to send to the receiver, we call this the sender's policy $s(\theta_S (i_L, i_R, t))\in V$
\item The receiver does not know the target, but sees the sender's symbol and tries to guess the target image. We call this the receiver's policy $r(i_L, i_R, s(\theta_S (i_L, i_R, t))) \in \lbrace L, R \rbrace$ 
\item If $r(i_L, i_R, s(\theta_S (i_L, i_R, t)) = t$, that is, if the receiver gets the target correct, both players receive a payoff of 1 (win), otherwise they receive a payoff of 0 (lose).
\end{enumerate}

Many extensions to the basic referential game explored in this paper are possible. There can be more images, a more sophisticated communication protocol (e.g., communication of a sequence of symbols or multi-step ``20 questions''-style communication requiring back-and-forth interaction),  rotation of the sender and receiver roles, having a human occasionally playing one of the roles, etc. 

Our setup is an instance of the cheap talk games studied in game theory%~\citep{crawford1982strategic}
. These games have been used as models to study the evolution of language in simulations in the lab~\citep{crawford1998survey,blume:1998experimental} and theoretically~\cite{crawford1982strategic}. They are interesting because they strip away all structural properties of language and focus on  its functional properties only. In cheap talk games, communication is just one symbol and the symbols have no ex-ante semantics: meaning only emerges in the context of the game.

Game theory is traditionally concerned with the study of Nash equilibria: stable policy profiles that no agent can deviate from to increase their payoffs. Cheap talk games with sufficient expressiveness and common interest admit a Nash equilibrium where all private information that the sender knows is revealed to the receiver~\citep{crawford1982strategic}. Thus, the question for us is whether learning agents will converge to this equilibrium (as convergence in learning is not guaranteed: \citealt{fudenberg2014recency}).  In addition, there are many equilibria which have the same amount of information revelation. However, some of the languages that emerge from these equilibria may not be the ones we are looking for, ie. they may capture high level semantic properties of our images. We tackle this issue empirically below, by analyzing the ``meanings'' that our agents assign to the symbols they learn to use. 

\section{Experimental Setup}
\label{sec:experimental-setup}
%We instantiate our referential game as follows.

\paragraph{Images} We use the McRae et al.'s
(\citeyear{McRae:etal:2005}) set of 463 base-level concrete
concepts (e.g., \emph{cat, apple, car}\ldots) spanning across 20
general categories (e.g., \emph{animal}, \emph{fruit/vegetable},
\emph{vehicle}\ldots). We randomly sample 100 images of each concept
from ImageNet~\citep{Deng:etal:2009}. To create target/distractor
pairs, we randomly sample two concepts, one image for each concept and
whether the first or second image will serve as target. We apply to
each image a forward-pass through the pre-trained VGG ConvNet~\cite{Simonyan:Zisserman:2014}, and represent it with the activations from either
the last softmax layer (\emph{sm}) or the second-to-last fully
connected layer (\emph{fc}).
%
% We start with a set of concepts studied by ~\cite{McRae:etal:2005}, in the context of property norm generation. This set contains 541 base-level concrete concepts (e.g., cat, apple, car etc.) that span across 20 general and broad categories (e.g., animal, fruit/vegetable, vehicle etc). For the purposes of the current experiments, \textbf{X} McRae concepts were excluded (either because of high ambiguity or for technical reasons), resulting in \textbf{Y} concepts.
% For each concept, we randomly sample 100 images from its respective ImageNet~\cite{Deng:etal:2009}  entry.  
%\COMMENT{This would be a good place to describe how the images are
%  represented by passing them through pre-trained CNNs, using one of
%  two layers}

\paragraph{Agent Players} Both sender and receiver are simple
feed-forward networks.  For the sender, we experiment with the two
architectures depicted in Figure \ref{fig:players}. Both take as input
the target and distractor representations with a pointer to the
target. The \emph{agnostic} sender is a generic neural network that
maps the original image vectors onto a ``game-specific'' embedding
space followed by a sigmoid activation function, and then applies fully-connected weights to their concatenation
to produce scores over vocabulary symbols.\COMMENT{Comment for future work: have you tried adding one more
  hidden layer above the embedding concatenation?}  The
\emph{informed} sender, just like the agnostic one, first embeds the
images into a ``game-specific'' space. It then applies 1-D
convolutions (``filters'') on the image embeddings by treating them as
different channels (i.e., convolutions with kernel size 2x1), followed by a sigmoid activation function (in
Figure~\ref{fig:players}, there are 4 such filters). The resulting
feature maps are combined through another filter (kernel size fx1, where f is the number of bottom filters),  to produce scores for
the vocabulary symbols. 
Intuitively, the informed sender compares the two image embeddings
dimension-by-dimension. For both senders, motivated by the discrete
nature of language, we enforce a strong
communication bottleneck that discretizes the communication
protocol. Activations  on the top (vocabulary) layer 
%is normalized by
%the softmax transform
are converted to a Gibbs distribution 
(with temperature parameter $\tau$), and then a
single symbol $s$ is sampled from the resulting probability
distribution.

The receiver takes as input the target and distractor image vectors in
randomized order and the symbol produced by the sender (as a one-hot
vector over the vocabulary).  It also embeds the images, as well as
the symbol, onto its own ``game-specific'' space. It then computes dot
products between the symbol and image embeddings. Ideally, dot similarity should be higher for the image that is better denoted by the symbol. %
% Intuitively, the dot similarity will be higher for the image for which
% the symbol embedding is most active. 
The two dot products are %softmax-transformed
converted to a  Gibbs distribution 
 (again with temperature parameter $\tau$) and the
receiver ``points'' to an image by sampling from the resulting
probability distribution.%\COMMENT{I changed the descriptions of
 % sampling of both sender and receiver to avoid a notation that is
%  different from the one used in the General Framework. If you want to
 % re-instated the notation, please make it compatible with the
 % notation in that section (original chunks commented out but still in
 % the source).}
% Finally, the receiver performs the action of ``pointing'' to
% an image $i$ by sampling from a Gibbs distribution
% $p_{\theta_R}(o|i_L,i_R,s)$ (with temperature $\tau$) parametrized
% by the scores over the two images $o_1$ and $o_2$ produced by the
% receiver.
% from a Gibbs distribution $p_{\theta_S}
% (s|o_t,o_d)$ (with temperature $\tau$) parametrized by the symbol
% scores produced by the sender. This produced symbol is the message
% that the sender sends over to the receiver.


%\COMMENT{Have you tried different depths for FULLYCONNECTED?} AL: Probably??
% I see I repeat myself :)

%\COMMENT{A bit more details: how are the networks trained, and vocabulary size}

\begin{figure}
\includegraphics[scale=0.20]{figures/players}
\caption{Architectures of agent players.%\textbf{Any comments on there or they look ok? Put a frame around the target for the sender architectures. You can perhaps also grey out the 3d symbol of the senders, to clarify the link with the third input to the receiver.}
}
\label{fig:players}
\end{figure}

\paragraph{General Training Details} We set the following
hyperparameters without tuning: embedding dimensionality of all embedding matrices: 50, number of filters applied to embeddings by
  informed sender: 20, temperature of all Gibbs distributions: 10. We explore two vocabulary sizes: \emph{10} and
\emph{100} symbols. The sender and receiver parameters $\theta=\langle
\theta_{R}, \theta_{S}\rangle$ are learned jointly while playing the
game. The only supervision used is communication success, i.e.,
whether the agents agreed on the referent. This setup is naturally
modeled with Reinforcement Learning \citep{Sutton:Barto:1998}. As
outlined in Section \ref{sec:general-framework}, the sender follows policy $s(\theta_S (i_L, i_R, t))\in V$ and the receiver policy $r(i_L, i_R, s(\theta_S (i_L, i_R, t))) \in \lbrace L, R \rbrace$. The loss function that the two agents must minimize is  $-\Expect_{\tilde{r}} [R(\tilde{r})]$ where $R$ is the reward function returning 1 iff $r(i_L, i_R, s(\theta_S (i_L, i_R, t)) = t$. %
%  The agent parameters implement a
% \emph{policy}. By executing it, the agents perform
% \emph{actions}, i.e., the sender picks a symbol from the distribution
% $p_{\theta_S} (s|o_t,o_d)$ and the receiver points to an image by
% sampling from the distribution $p_{\theta_R} (o|o_1,o_2,s)$. The loss
% function that the two agents are minimizing is
% $-\Expect_{\tilde{o}\sim p(o|o_1,o_2,s)}[R(\tilde{o})]$, where $R$ is
% the reward function which returns 1 iff $\tilde{o} =$ target. 
Parameter are updated through the Reinforce rule~\cite{Williams:1992}.
We apply mini-batch updates, with a batch size of 32 and for a total of
50k iterations (games).  \textbf{At test time, we compile a set of 10k games  consisting of a target and a distractor by randomly sampling two objects and for each object an image.}

We now turn to our main questions. The first is whether the agents can
learn to successfully coordinate in a reasonable amount of time (as
there are plenty of examples where certain learning models are simply
unable to converge to Nash equilibria, given certain inputs
\textbf{(CITE)}). The second is whether the language that agents
converge to can be thought of as ``natural language'', in the sense
that symbols are assigned meanings that make intuitive sense in terms
of our conceptualization of the world.








